\documentclass[11pt]{article}
\usepackage{mathblatt}
% Präambel des Setzers (werkzeuge/setzer-praeambel.tex, wortgleich); dazu \gteil (Muster T2U)
\makeatletter
\setlength{\mbpfbe}{0mm}% keine Punkte (pruefung.md § 5: nur beim Prüfstein)
\newlength{\gEin}\setlength{\gEin}{0pt}
\let\gSchrift\relax
\renewcommand{\pfkreuzzeile}[1]{\par\noindent\hangindent3.2em\hangafter1\hspace*{1.6em}$\square$\ #1\par}
\newcommand{\gkreuz}[1]{$\square$\ #1\hspace{2em}}
\newcommand{\gkreuzzeile}[2]{\par\noindent\hangindent2.2em\hangafter1\mbox{#1\,$\square$}\ #2\par}% Buchstabe, Option: eine je Zeile (neben Skizze)
% links, was man ansieht, rechts, was man tut; Spaltengrenze überall an derselben Stelle (gemeinsam.md „Satz“)
% v0.17: Fußeintrag „Nr: Wert“ – kurz in einer Box (nie getrennt), lang umbrechbar mit der Nummer am Anfang
% (vorher \mbox: lange Lösungen liefen über den Rand, AS40 138 pt); robust, weil \fusshilfe in die .aux schreibt
\NewDocumentCommand{\gfuss}{m m}{\sbox0{#1:~#2}\ifdim\wd0<0.45\textwidth\mbox{#1:~#2}\else\mbox{#1:}\nobreak\ #2\fi}
\newsavebox{\gbox}\newlength{\gLinks}
% v0.10 (Lehrer 13.10.): keine starre Halbierung – die linke Spalte ist so breit wie die Figur (+ 6 mm), die
% Antwort steht direkt daneben; ist die Figur breiter als 0,62 Zeile, steht die Antwort darunter
\newcommand{\gzwei}[2]{\par\smallskip\noindent\sbox{\gbox}{#1}%
  \setlength{\gLinks}{\dimexpr\wd\gbox+6mm\relax}%
  \ifdim\gLinks>0.62\linewidth
    \usebox{\gbox}\par\smallskip\noindent\begin{minipage}{\linewidth}#2\end{minipage}%
  \else
    \begin{minipage}[t]{\gLinks}\vspace{0pt}\usebox{\gbox}\end{minipage}%
    \begin{minipage}[t]{\dimexpr\linewidth-\gLinks\relax}\vspace{0pt}\raggedright #2\end{minipage}%
  \fi\par}
% 16.10. (Lehrer, PA10 S. 1): Text links, Koordinatensystem rechts daneben; ist das System breiter als 0,58 Zeile,
% steht es wie bisher unter dem Text
% Passt das System knapp nicht mehr auf die Seite (bis 25 % zu hoch), wird es auf den Rest verkleinert statt
% die Seite halb leer zu lassen
\newlength{\gRest}\newlength{\gHoch}
\newcommand{\gtextks}[2]{\sbox{\gbox}{#1}%
  \setlength{\gHoch}{\dimexpr\ht\gbox+\dp\gbox\relax}%
  \ifdim\pagegoal<\maxdimen\setlength{\gRest}{\dimexpr\pagegoal-\pagetotal-4\baselineskip\relax}%
    \ifdim\gHoch>\gRest\ifdim\gRest>0.75\gHoch
      \sbox{\gbox}{\resizebox*{!}{\gRest}{\usebox{\gbox}}}\fi\fi\fi
  \ifdim\wd\gbox>0.58\linewidth #2\par\smallskip\noindent\gpassend{\usebox{\gbox}}\par
  \else\begin{minipage}[t]{\dimexpr\linewidth-\wd\gbox-5mm\relax}\raggedright #2\end{minipage}\hfill
    \begin{minipage}[t]{\wd\gbox}\vspace{0pt}\usebox{\gbox}\end{minipage}\par\fi}
% v0.9 (12.10., pruefung.md § 5 „Marke mit B, Bedingungen rechts“): drittes Argument von pfaufg = Bedingungen
% der Prüfung (v0.10: \pfohnetr = TR durchgestrichen, „FOR“), klein grau im rechten Rand auf der Grundlinie der Nummer, Zeilen mit \\
\newcommand{\gbedrand}[2]{\ifblank{#2}{}{\rlap{\hspace*{\dimexpr#1+2mm\relax}%
  \parbox[t]{15mm}{\raggedright\scriptsize\leavevmode\color{mbgrau}\strut#2}}}}
% Teil eines Bündels (a, b …): Marke links im Rand, Bedingungen rechts im Rand, auf der Zeile des Teils
\newcommand{\gteilrand}[2]{\llap{\mbpfrandmarke{#1}\hspace*{\mbpfnr}}\gbedrand{\linewidth}{#2}}
% Teil eines Bündels: Buchstabe in eigener Spalte, Inhalt eingerückt (\gEin = Einzug, damit die Spaltengrenze von
% \gzwei an derselben Stelle bleibt)
\newenvironment{gbteil}[3]{\par\smallskip\noindent\gteilrand{#2}{#3}\makebox[\mbpfteil][l]{\textbf{#1}}%
  \begin{minipage}[t]{\dimexpr\linewidth-\mbpfteil\relax}\setlength{\gEin}{\mbpfteil}\raggedright\strut\ignorespaces}%
 {\end{minipage}\par\addvspace{6pt}}% v0.18 (PG80 Nr. 1): Abstand, damit die Linie von a) nicht an b) rutscht
% Grafik höchstens so breit wie die Spalte (Bündel nebeneinander, Koordinatensystem ohne Karo)
\newsavebox{\gpbox}
\newcommand{\gpassend}[1]{\sbox{\gpbox}{#1}\ifdim\wd\gpbox>\linewidth\typeout{GPASSEND: verkleinert \the\wd\gpbox\space auf \the\linewidth}\resizebox{\linewidth}{!}{\usebox{\gpbox}}\else\usebox{\gpbox}\fi}
% Skizzenfeld (v0.10, Lehrer 13.10.): Karo ab dem linken Rand über die Zeile, „Skizze“ klein grau oben links im Feld
% (wie \pfskizzenfeld der Vorlage); der rechte Rand bleibt den Bedingungen
\newcommand{\gskizze}[1]{\par\smallskip\noindent\begin{tikzpicture}
  \clip (0,0) rectangle ({\linewidth-1mm},{#1*10mm});
  \draw[black!16,line width=0.3pt,step=5mm] (0,0) grid ({\linewidth},{#1*10mm});
  \node[anchor=north west,inner sep=2pt,font=\scriptsize,text=mbgrau,fill=white] at (0,{#1*10mm}) {Skizze};
  \end{tikzpicture}\par}
\RenewDocumentEnvironment{pfaufg}{m m m}{%
  \begin{lrbox}{\mbaufgabebox}\begin{minipage}[t]{\linewidth}%
  \gSchrift\noindent\hspace*{\gEin}\mbpfrandmarke{#1}\gbedrand{\linewidth-\gEin}{#3}%
  \begin{minipage}[t]{\mbpfnr}\strut\textbf{#2}\end{minipage}%
  \begin{minipage}[t]{\dimexpr\linewidth-\gEin-\mbpfnr\relax}\raggedright\strut\ignorespaces}%
 {\end{minipage}%
  \end{minipage}\end{lrbox}%
  \setlength{\mbaufgabehoehe}{\dimexpr\ht\mbaufgabebox+\dp\mbaufgabebox\relax}%
  \ifdefempty{\gkopfnach}{\par\addvspace{5pt}\Needspace*{\mbaufgabehoehe}}{%
    % Stufenkopf, der mit dieser (ersten) Aufgabe zusammenbleibt: Platz für beide verlangen, dann den Kopf setzen
    % (12.10., Quadratische: ein Kopf stand allein unten, die Aufgabe mit Koordinatensystem auf der nächsten Seite)
    % v0.17 (AS40 „Terme deuten“ allein unten): hinter dem Kopf steht das Sprungziel (\special, ein Knoten, nach
    % dem die Seite brechen durfte) – darum \nobreak nach dem Kopf; Vorhalt 56 pt = Kopf mit Linie und Abständen
    \par\Needspace*{\dimexpr\mbaufgabehoehe+56pt\relax}\gkopfnach\par\nobreak\gdef\gkopfnach{}\addvspace{5pt}}%
  \noindent\usebox{\mbaufgabebox}\par\addvspace{5pt}}
% v0.18 (plan.md § 10, 15.10. l): Abschnittsbuchstabe A, B, C vor dem Namen im Kopf – als Vorsatz, damit der
% Name im TeX gleich bleibt (Bestellkarte, Sprungziele lesen ihn)
\def\gbstakt{}\newcommand{\gbstvor}[1]{\gdef\gbstakt{#1}}
\NewCommandCopy\gpfstufekopfalt\pfstufekopf
\RenewDocumentCommand{\pfstufekopf}{o m m}{\IfValueTF{#1}{\gpfstufekopfalt[#1]}{\gpfstufekopfalt}%
  {\ifx\gbstakt\empty\else\gbstakt\enspace\fi #2}{#3}\gdef\gbstakt{}}
% v0.18 (16.10. e): Begriff grau, ohne Nummer, ohne Rahmen über der ersten Aufgabe der Stufe
\newcommand{\gbegriff}[1]{\par\nobreak\noindent\hspace*{\mbpfnr}\parbox[t]{\dimexpr\linewidth-\mbpfnr\relax}{\raggedright\small\color{mbgrau}#1}\par\nobreak\addvspace{3pt}}
% v0.18 (16.10. g): Lösungsblatt – Kopf mit kräftigem Strich, Abschnitt klein grau als Zwischenzeile
\newcommand{\glkopf}[1]{\noindent{\large\bfseries Lösungen \textperiodcentered{} #1}\par\vspace{3pt}\hrule height 1.2pt\par\medskip}
\newcommand{\glzwischen}[1]{\par\addvspace{8pt}\Needspace*{4\baselineskip}\noindent{\footnotesize\color{mbgrau}#1}\par\nobreak\vspace{1pt}}
% v0.18 (15.10. o, w): Blattende – Kern ohne Rahmen, dann graue Zeile „Mehr: n Aufgaben“
\newcommand{\gkern}[1]{\par\addvspace{14pt}\Needspace*{6\baselineskip}\noindent{\bfseries Kern}\par\nobreak\smallskip
  {\small\raggedright #1\par}}
\newcommand{\gmehr}[1]{\par\addvspace{10pt}\noindent{\footnotesize\color{mbgrau}Mehr: #1 Aufgaben}\par}
% Stufenkopf vormerken: \gkopfvor{<Kopf>} – die nächste pfaufg setzt ihn, sobald beide zusammen Platz haben
\def\gkopfnach{}
\newcommand{\gkopfvor}[1]{\gdef\gkopfnach{#1}}
% Teilaufgabe mit Skizze links in gzwei: Buchstabe in eigener Spalte, darunter Text und Skizze
\newcommand{\gteil}[2]{\makebox[\mbpfteil][l]{\textbf{#1}}\parbox[t]{\dimexpr0.5\textwidth-\gEin-\mbpfnr-5mm-\mbpfteil\relax}{\raggedright #2}}
% Stufenkopf der Vorlage (mathblatt.sty 2026-10-12: Linie, Abstand, 13 pt, \pfjahre) ohne ihren Vorhalt von
% 14 Zeilen: der Setzer fordert vor dem Kopf 0,16 Seitenhöhe (Kopf + erste Aufgabe), mehr nicht
\setlength{\mbpfstufevorhalt}{0pt}
\makeatother
\newcommand{\karo}[1]{\par\smallskip\noindent\begin{tikzpicture}
  \clip (0,0) rectangle ({\linewidth-1mm},{#1*10mm});
  \draw[black!16,line width=0.3pt,step=5mm] (0,0) grid ({\linewidth},{#1*10mm});
  \end{tikzpicture}\par}
\newcommand{\kk}[1]{$\square$\,#1\hspace{0.9em}}
\newcommand{\linie}[1]{\,{\color{black!45}\rule[-1pt]{#1}{0.4pt}}\,}
\newcommand{\pkt}{\textperiodcentered{}}
% v0.14 (14.10.): Stufenkopf mit fünf Punkten (Option = Jahre für die Bestellkarte) oder „zuletzt 19“
\NewDocumentCommand{\gpunkte}{o m}{#2}
\newcommand{\gzuletzt}[1]{zuletzt #1}
% v0.7 (12.10.): Schülerfrage unter dem Stufenkopf; Ankreuzoption als Bild (Buchstabe und Kästchen darunter)
\newcommand{\gfrage}[1]{\par\nobreak\noindent\hspace*{\mbpfnr}{\small\itshape #1}\par\nobreak\addvspace{2pt}}
\newlength{\gbildbreite}\setlength{\gbildbreite}{0.32\linewidth}% Setzer setzt sie je Reihe (vier Bilder: 0,24)
\newcommand{\gbildoption}[3]{\begin{minipage}[t][#3][b]{\gbildbreite}\centering #2\par\smallskip\mbox{#1\,$\square$}\end{minipage}}
% Lineare (12.10.): Bild mit Buchstabe ohne Kästchen (Diagramm zuordnen); Aussage mit „wahr / falsch“ zum Ankreuzen
\newcommand{\gbildname}[3]{\begin{minipage}[t][#3][b]{\gbildbreite}\centering #2\par\smallskip\mbox{\textbf{#1}}\end{minipage}}
\newcommand{\gwahrfalsch}[2]{\par\noindent\hangindent2.2em\hangafter1\makebox[2.2em][l]{(#1)}#2\hfill\mbox{wahr $\square$\quad falsch $\square$}\par}

\newcommand{\gunter}[2]{\par\addvspace{6pt}\noindent{\bfseries #1\enspace #2}\par\nobreak\vspace{2pt}}
\newcommand{\gsternzeile}{\par\vspace{-4pt}\noindent{\footnotesize\color{mbgrau}* nur FOR}\par\medskip}
% 17.10.: Skizze oben bündig mit der ersten Textzeile (beide Spalten hängen an der Oberkante); der Rechenraum (Karo)
% beginnt direkt unter dem Text, links neben dem Graph, und läuft unter dem Graph über die ganze Zeile weiter
% (Streifen links = auf 5 mm aufgerundete Höhe des Graphen unter dem Text, Raster durchgehend; die Fläche des Streifens
% wird von der Karohöhe unter dem Graph abgezogen, auf 0,5 cm aufgerundet – die Gesamtfläche Rechenraum bleibt mindestens #3 × Zeile); #3 = Zeilen Karo in cm
% unter dem Graph (0 = kein Karo). Kein Verkleinern des Graphen.
\newsavebox{\gtbox}\newlength{\gWl}\newlength{\gDh}
\newcommand{\gkaroohne}[1]{\nointerlineskip\noindent\begin{tikzpicture}
  \clip (0,0) rectangle ({\linewidth-1mm},{#1*10mm});
  \draw[black!16,line width=0.3pt,step=5mm] (0,0) grid ({\linewidth},{#1*10mm});
  \end{tikzpicture}\par}
\newcommand{\gtextkaro}[3]{\sbox{\gbox}{#1}%
  \ifdim\wd\gbox>0.58\linewidth #2\par\smallskip\noindent\gpassend{\usebox{\gbox}}\par\ifnum#3>0 \karo{#3}\fi
  \else
    \setlength{\gWl}{\dimexpr\linewidth-\wd\gbox-5mm\relax}%
    \sbox{\gtbox}{\begin{minipage}[t]{\gWl}\raggedright #2\par\end{minipage}}%
    \setlength{\gDh}{\dimexpr\ht\gbox+\dp\gbox-\ht\gtbox-\dp\gtbox-6pt\relax}%
    \leavevmode
    \ifnum#3>0\ifdim\gDh>0pt
      \pgfmathsetmacro{\gxs}{floor(\gWl/5mm)*5mm}\pgfmathsetmacro{\gys}{ceil(\gDh/5mm)*5mm}%
      \begin{minipage}[t]{\gWl}\usebox{\gtbox}\par\vspace{6pt}\nointerlineskip\noindent
        \begin{tikzpicture}
        \clip (0,0) rectangle (\gxs pt,\gys pt);
        \draw[black!16,line width=0.3pt,step=5mm] (0,0) grid (\gxs pt,\gys pt);
        \end{tikzpicture}\par\end{minipage}%
    \else\usebox{\gtbox}\fi\else\usebox{\gtbox}\fi
    \hfill\begin{minipage}[t]{\wd\gbox}\vspace{-\dimexpr\ht\strutbox+2.4ex\relax}\usebox{\gbox}\end{minipage}\par
    \ifnum#3>0 \ifdim\gDh>0pt
      \pgfmathsetmacro{\gnn}{max(0,ceil(2*(#3-(\gys/28.4528)*(\gxs/\linewidth)))/2)}%
      \ifdim\gnn pt>0pt\gkaroohne{\gnn}\fi
    \else\karo{#3}\fi\fi
  \fi}
\begin{document}

\pfheftstil
\pflaufkopf[P10 \pkt{} Prüfung]{Gleichungen und Gleichungssysteme}{PA20}\pfseitevon
\pfheftkopf{Gleichungen und Gleichungssysteme}{}
\pfstufenzeile{Lineare Gleichungen \pkt{} Lineare Gleichungssysteme}
\gsternzeile

% ===================== A Lineare Gleichungen =====================
\gkopfvor{\gbstvor{A}\pfstufekopf{Lineare Gleichungen}{}}
% DRUCK 2024-OS-B1d deckt 2024-OS-B1d
\begin{pfaufg}{nach~P10\\’24}{1.}{Basis}
Gib die Lösung der Gleichung $2\cdot(x-6{,}5)=0$ an.
\karo{3}
\par\medskip\noindent\hfill x $=$\linie{25mm}
\fusshilfe{\gfuss{1}{6,5}}
\end{pfaufg}
% DRUCK 2020-OS-B1e deckt 2020-OS-B1e
\begin{pfaufg}{nach~P10\\’20}{2.}{Basis}
Gib die Lösung der Gleichung $2(x-4)=6$ an.
\karo{3}
\par\medskip\noindent\hfill x $=$\linie{25mm}
\fusshilfe{\gfuss{2}{7}}
\end{pfaufg}
% DRUCK 2023-OS-B1e deckt 2023-OS-B1e
\begin{pfaufg}{nach~P10\\’23}{3.}{Basis}
Gib die Lösung der Gleichung $3\cdot(x-8)+2=2$ an.
\karo{4}
\par\medskip\noindent\hfill x $=$\linie{25mm}
\fusshilfe{\gfuss{3}{8}}
\end{pfaufg}
% DRUCK 2018-OS-B1c deckt 2018-OS-B1c
\begin{pfaufg}{nach~P10\\’18}{4.}{Basis}
Kreuze an, für welche Gleichung $x=-2$ gilt:\pfkreuzab{}{$4x-8=0$}\pfkreuzab{}{$2x+10=2$}\pfkreuzab{}{$5x+12=2$}\pfkreuzab{}{$-2x+4=0$}
\fusshilfe{\gfuss{4}{$5x+12=2$}}
\end{pfaufg}
% DRUCK 2022-OS-B1c deckt 2022-OS-B1c
\begin{pfaufg}{nach~P10\\’22}{5.}{Basis}
Gegeben ist die Gleichung $5-2x=3x-25$. Kreuze an, welche Zahl die Lösung der Gleichung ist:\pfkreuzab{}{$x=2$}\pfkreuzab{}{$x=4$}\pfkreuzab{}{$x=6$}\pfkreuzab{}{$x=8$}
\karo{3}
\fusshilfe{\gfuss{5}{$x=6$}}
\end{pfaufg}
% DRUCK 2021-OS-B1e deckt 2021-OS-B1e
\begin{pfaufg}{nach~P10\\’21}{6.}{Basis}
Das Fünffache einer Zahl vermindert um 4 ist gleich 36. Kreuze die richtige Gleichung an:\pfkreuzab{}{$5x=-4+36$}\pfkreuzab{}{$5-4x=36$}\pfkreuzab{}{$5x-4=36$}\pfkreuzab{}{$36-4=5x$}
\fusshilfe{\gfuss{6}{$5x-4=36$}}
\end{pfaufg}
% DRUCK 2016-OS-B1b deckt 2016-OS-B1b
\begin{pfaufg}{nach~P10\\’16}{7.}{Basis}
Stelle zu folgender Formulierung die passende Gleichung auf: Das Achtfache einer Zahl vermindert um zwölf ist gleich 36.
\par\medskip\noindent\hfill \linie{60mm}
\fusshilfe{\gfuss{7}{36}}
\end{pfaufg}
% DRUCK 2018-OS-K3b deckt 2018-OS-K3b
\begin{pfaufg}{nach~P10\\’18}{8.}{}
\gtextkaro{{\small\renewcommand{\arraystretch}{1.15}\begin{tabular}{|l|r|}\hline Lieblingsessen & Anzahl\\\hline Schnitzel & 80\\ Pizza & 90\\ Kartoffelgerichte & 160\\ Gemüsegerichte & 160\\ Fischgerichte & 110\\ Pasta & ?\\ Salat & ?\\\hline\end{tabular}}}{1000 Jugendliche wurden nach ihrem Lieblingsessen befragt. In der Tabelle fehlen zwei Angaben. Bekannt ist: Die Anzahl der Jugendlichen, die Pasta als Lieblingsessen angegeben hat, ist dreimal höher als die Anzahl der Jugendlichen, die Salat als Lieblingsessen bevorzugt. Ermittle die fehlenden Angaben.}{5}
\par\medskip\noindent\hfill Pasta $=$\linie{25mm}\quad Salat $=$\linie{25mm}
\fusshilfe{\gfuss{8}{300; 100}}
\end{pfaufg}
% DRUCK 2020-OS-K2e deckt 2020-OS-K2e
\begin{pfaufg}{nach~P10\\’20}{9.\textsuperscript{*}}{}
Am Wochenende kamen 520 Besucher in das Schwimmbad. An diesen Tagen kostet die Tageskarte 12 € pro Person. Der Betreiber des Schwimmbades behauptet: „Wären am Wochenende $x$ Besucher mehr gekommen, dann hätte ich am Wochenende 7920 € Einnahmen erzielt.“ Stelle eine Gleichung auf, die diesen Sachverhalt beschreibt. Ermittle die Anzahl $x$ der Besucher, die dafür gefehlt hat.
\karo{5}
\par\medskip\noindent\hfill x $=$\linie{25mm}
\fusshilfe{\gfuss{9}{140}}
\end{pfaufg}

% ===================== B Lineare Gleichungssysteme =====================
\gkopfvor{\gbstvor{B}\pfstufekopf{Lineare Gleichungssysteme}{}}
% DRUCK 2022-OS-K7a deckt 2022-OS-K7a
% DRUCK 2022-OS-K7b deckt 2022-OS-K7b
\begin{pfaufg}{nach~P10\\’22}{10.}{}
Frau Bauer kauft für ihren Garten Apfelbäume und Pflaumenbäume. Ein Apfelbaum und ein Pflaumenbaum kosten zusammen 63,00 €. Für 9 Apfelbäume und 3 Pflaumenbäume zahlt sie 352,20 €. Frau Bauer stellt dazu zwei Gleichungen auf: I $x + y = 63$ und II $9x + 3y = 352{,}20$.
\begin{gbteil}{a)}{}{}
Gib die Bedeutung der Variablen $x$ und $y$ an.
\par\medskip\noindent\hfill x $=$\linie{45mm}
\par\medskip\noindent\hfill y $=$\linie{45mm}
\end{gbteil}
\begin{gbteil}{b)}{}{}
Berechne den Preis für einen Apfelbaum und den Preis für einen Pflaumenbaum.
\karo{6}
\par\medskip\noindent\hfill Apfelbaum $=$\linie{25mm}\quad Pflaumenbaum $=$\linie{25mm}
\end{gbteil}
\fusshilfe{\gfuss{10}{a) Preis Apfelbaum in €; Preis Pflaumenbaum in €; b) 27,20 €; 35,80 €}}
\end{pfaufg}\vspace{-6pt}
% DRUCK 2024-OS-K7a deckt 2024-OS-K7a
% DRUCK 2024-OS-K7b deckt 2024-OS-K7b
\begin{pfaufg}{nach~P10\\’24}{11.}{}
Die Geschwister Daniel und Filip wollen ihrer Mutter jeweils einen Blumenstrauß zum Muttertag kaufen. Daniel kauft seinen Blumenstrauß in Geschäft A, Filip kauft seinen in Geschäft B.
\begin{gbteil}{a)}{}{}
Im Geschäft A kosten eine Rose und eine Tulpe zusammen 3,80 €. Für 6 Rosen und 5 Tulpen zahlt Daniel insgesamt 21,20 €. Stelle ein lineares Gleichungssystem (I und II) zu diesem Sachverhalt auf. Dabei gilt: $x$: Preis pro Rose, $y$: Preis pro Tulpe.
\par\medskip\noindent\hfill \linie{60mm}
\par\medskip\noindent\hfill \linie{60mm}
\end{gbteil}
\begin{gbteil}{b)}{}{}
Im Geschäft B ist der Preis für eine Rose 2,30 € und der Preis für eine Tulpe 1,70 €. Filip stellt zu seinem Blumenstrauß dieses Gleichungssystem auf: I. $2{,}30\,r + 1{,}70\,t = 26{,}90$ und II. $r + t = 13$. Formuliere einen Satz, der den Sachverhalt beschreibt, der durch Gleichung II ausgedrückt wird. Ermittle die Lösung für das Gleichungssystem.
\karo{6}
\par\medskip\noindent\hfill \linie{60mm}
\par\medskip\noindent\hfill r $=$\linie{25mm}\quad t $=$\linie{25mm}
\end{gbteil}
\fusshilfe{\gfuss{11}{a) 3,80; 21,20; b) 13 Blumen im Strauß; 8; 5}}
\end{pfaufg}\vspace{-6pt}
% DRUCK 2021-OS-K7b deckt 2021-OS-K7b
\begin{pfaufg}{nach~P10\\’21}{12.}{}
Zwei Familien besuchen in Paris den Eiffelturm. Familie Beyer (zwei Erwachsene und zwei Kinder) zahlt 77,80 € Eintritt. Familie Gouvan (ein Erwachsener und drei Kinder) zahlt 64,90 € Eintritt. Ermittle den Eintrittspreis für Erwachsene und den Eintrittspreis für Kinder. Stelle zunächst zwei Gleichungen zu diesem Sachverhalt auf.
\karo{7}
\par\medskip\noindent\hfill I. $=$\linie{45mm}
\par\medskip\noindent\hfill II. $=$\linie{45mm}
\par\medskip\noindent\hfill Erwachsene $=$\linie{25mm}\quad Kinder $=$\linie{25mm}
\fusshilfe{\gfuss{12}{77,80; 64,90; 25,90 €; 13,00 €}}
\end{pfaufg}\vspace{-6pt}
% DRUCK 2016-OS-K6d deckt 2016-OS-K6d
\begin{pfaufg}{nach~P10\\’16}{13.\textsuperscript{*}}{}
In der Jugendherberge gibt es Drei-Bett-Zimmer und Fünf-Bett-Zimmer. Es stehen 16 Zimmer mit insgesamt 66 Betten zur Verfügung. Ermittle die Anzahl der Drei-Bett-Zimmer und der Fünf-Bett-Zimmer.
\karo{7}
\par\medskip\noindent\hfill Drei-Bett-Zimmer $=$\linie{25mm}\quad Fünf-Bett-Zimmer $=$\linie{25mm}
\fusshilfe{\gfuss{13}{7; 9}}
\end{pfaufg}\vspace{-6pt}
\end{document}